\documentclass[四库全书文渊阁简明目录]{ltc-guji}

\usepackage{enumitem} % For better list control if needed
\usepackage{tikz}
% \禁用分页裁剪
\无标点模式
\setmainfont{TW-Kai}


\title{欽定四庫全書簡明目錄}

\chapter{經部 易類\\卷一}

\begin{document}
\begin{正文}
欽定四庫全書簡明目錄卷一

\条目[1]{經部一}
\条目[2]{易類}

子夏易傳十一卷

\注{舊本題卜子夏撰，實後人輾轉依托，非其原書。然唐、宋以來，流傳已久，今仍錄冠易類之首。凡托名之書，仍從其所托之時代，《漢書·藝文志》例也。}

\按{謹按：唐徐堅《學記》以太宗御制升列歷代之前，蓋尊尊之大義。焦竑《國史經籍志》，朱彞尊《經義考》並踵前規。臣等編摩《四庫》，初亦恭錄

《御定易經通注》\\
《御纂周易折中》\\
《御纂周易述義》，弁冕諸經。仰蒙\\
指示，命冠於

國朝著述之首，俾尊卑有序，而時代不淆。

聖度謙沖，酌中立憲，實為千古之大公。謹恪遵彞訓，仍託始於《子夏易傳》。並發凡於此，著《四庫》之通例焉。}

\end{正文}
\end{document}
